% =====================================================================
% Slide 4: Young's modulus E for mechanical engineers
% =====================================================================
\begin{frame}{Young's Modulus $E$ --- The Mechanical Engineer's View}
\begin{columns}[T]
\column{0.55\textwidth}
\textbf{Definition.} In uniaxial loading,
$E = \dfrac{\sigma}{\varepsilon}$ --- axial stress per unit axial
strain (units: GPa).

\medskip
\textbf{Where $E$ appears in design:}
\begin{itemize}\small
  \item \textbf{Beam deflection.} Cantilever tip deflection
        $\delta = \dfrac{F L^3}{3 E I}$ --- a $2\times$ change in
        $E$ doubles tip rigidity for the same beam.
  \item \textbf{Buckling load.} Euler $P_{\mathrm{cr}} =
        \dfrac{\pi^2 E I}{(KL)^2}$ --- compression-member capacity
        scales linearly with $E$.
  \item \textbf{Natural frequency.} For a clamped beam
        $\omega_n \propto \sqrt{E/\rho}$ --- vibration design,
        modal analysis, NVH (noise-vibration-harshness).
  \item \textbf{Hertz contact stress.} Bearing/gear contact area
        and pressure depend on the effective $E^{*}$ of mating
        bodies.
\end{itemize}

\column{0.43\textwidth}
\begin{block}{Typical engineering ranges}
\small
\begin{tabular}{lr}
\toprule
Material           & $E$ (GPa) \\
\midrule
Polymers (PE, PP)  & 0.5 -- 3 \\
Magnesium alloys   & 40 -- 45 \\
\textbf{Aluminium 7075} & \textbf{70 -- 75} \\
Titanium 6Al-4V    & 110 \\
Steel (mild)       & 200 \\
Tungsten           & 410 \\
\bottomrule
\end{tabular}
\end{block}
\vspace{0.3em}
\small\textit{$E$ is what designers minimise when they want
flexibility (springs, panels) and maximise when they want rigidity
(machine frames, beams).}
\end{columns}
\end{frame}

% =====================================================================
% Slide 5: Poisson's ratio nu for mechanical engineers
% =====================================================================
\begin{frame}{Poisson's Ratio $\nu$ --- The Mechanical Engineer's View}
\begin{columns}[T]
\column{0.55\textwidth}
\textbf{Definition.}
$\nu = -\dfrac{\varepsilon_{\text{lateral}}}{\varepsilon_{\text{axial}}}$ ---
the lateral contraction per unit axial extension (dimensionless).

\medskip
\textbf{Where $\nu$ appears in design:}
\begin{itemize}\small
  \item \textbf{Plane-strain vs plane-stress.} The transition
        criterion in thin- vs thick-walled pressure vessels and
        sheet-metal forming hinges on $\nu$.
  \item \textbf{Volumetric vs deviatoric coupling.} Bulk modulus
        $K = \dfrac{E}{3(1-2\nu)}$ diverges as $\nu \!\to\! 0.5$ ---
        incompressible limit. Critical for rubber seals, soft
        polymers, energy-absorbing foams (auxetic if $\nu<0$).
  \item \textbf{Stress concentration.} Notch and hole stress
        concentration factors at finite thickness depend on $\nu$
        through the Lam\'e parameters.
  \item \textbf{Fatigue \& fracture.} Plane-strain $K_{Ic}$ vs
        plane-stress $K_c$ thresholds differ by a $(1-\nu^2)$
        factor.
\end{itemize}

\column{0.43\textwidth}
\begin{block}{Stability bounds (isotropic)}
\small
$-1 < \nu < 0.5$.
\\[0.3em]
For \emph{cubic} crystals, the Born--Cauchy criterion
$C_{11} > C_{12}$ is the equivalent mechanical-stability bound;
violating it gives a negative shear modulus.
\end{block}
\begin{block}{Typical ranges}
\small
\begin{tabular}{lr}
\toprule
Material         & $\nu$ \\
\midrule
Cork (auxetic-ish) & $\sim\!0.0$ \\
\textbf{Al alloys} & \textbf{0.33} \\
Steel              & 0.30 \\
Rubber             & $\to\!0.5$ \\
Beryllium          & 0.03 \\
\bottomrule
\end{tabular}
\end{block}
\end{columns}
\end{frame}

% =====================================================================
% Slide 6: From C11, C12 to (E, nu) and stability
% =====================================================================
\begin{frame}{Why We Train on $C_{11}, C_{12}$, Not on $(E,\nu)$ Directly}
\begin{columns}[T]
\column{0.55\textwidth}
For a cubic crystal the engineering moduli are recovered
analytically from the two independent stiffness constants:
\begin{equation*}
  E = \frac{(C_{11}-C_{12})(C_{11}+2C_{12})}{C_{11}+C_{12}},
  \qquad
  \nu = \frac{C_{12}}{C_{11}+C_{12}}.
\end{equation*}

\vspace{0.4em}
\textbf{The pitfall.}\;
$\nu = C_{12}/(C_{11}+C_{12})$ is ill-conditioned ---
a $5\,\%$ error on each $C_{ij}$ can blow up to a
$30\,\%$ relative error on $\nu$ for Mg/Zn-rich compositions
where the denominator is small.

\vspace{0.5em}
\textbf{Decision.} Both ML surrogates (composition NN and the
NN-MEAM optimizer) train under Huber loss
\emph{on $C_{11}$ and $C_{12}$}, and recover $(E,\nu)$
\emph{after} prediction.

\column{0.43\textwidth}
\begin{block}{Physicality filter (applied everywhere)}
\small
A configuration is dropped if any of:
\begin{itemize}\scriptsize
  \item $E < 0$
  \item $\nu < 0$
  \item $\nu \ge 0.48$ (guard for $1\!-\!2\nu$ singularity)
  \item $C_{11} < C_{12}$ (Born--Cauchy violation)
\end{itemize}
The same filter is enforced in Stage 2 (LAMMPS),
Stage 5a (composition NN), and Stage 5b (NN-MEAM optimizer)
to keep training, validation, and verification on the same playing
field.
\end{block}
\end{columns}
\end{frame}
